\documentclass[UTF8,11pt,fontset=windows]{ctexart}

\usepackage[a4paper,top=2.4cm,bottom=2.3cm,left=2.6cm,right=2.6cm]{geometry}
\usepackage{amsmath}
\usepackage{booktabs}
\usepackage{tabularx}
\usepackage{array}
\usepackage{enumitem}
\usepackage{xcolor}
\usepackage{hyperref}
\usepackage{fancyhdr}
\usepackage{microtype}

\definecolor{accent}{HTML}{234E70}
\definecolor{lightline}{HTML}{D7DEE5}
\hypersetup{
  colorlinks=true,
  linkcolor=accent,
  citecolor=accent,
  urlcolor=accent,
  pdftitle={Schema-Aware Candidate SQL Selection: Project Progress Report},
  pdfauthor={Dongx Xia}
}

\setlength{\parindent}{2em}
\setlength{\parskip}{0.25em}
\setlist[itemize]{leftmargin=2em,itemsep=0.2em,topsep=0.25em}
\setlist[enumerate]{leftmargin=2.2em,itemsep=0.2em,topsep=0.25em}
\renewcommand{\arraystretch}{1.18}
\pagestyle{fancy}
\fancyhf{}
\fancyhead[L]{\small Schema-Aware Candidate SQL Selection}
\fancyhead[R]{\small Dongx Xia}
\fancyfoot[C]{\thepage}
\renewcommand{\headrulewidth}{0.4pt}

\ctexset{
  section={format=\Large\bfseries\color{accent},beforeskip=1.0em,afterskip=0.45em},
  subsection={format=\large\bfseries,beforeskip=0.75em,afterskip=0.3em}
}

\title{\vspace{-1.8em}\textbf{Schema-Aware Candidate SQL Selection}\\[0.35em]
\large 项目阶段性汇报}
\author{Dongx Xia \quad Student ID: 50062682}
\date{2026年9月}

\begin{document}
\maketitle
\vspace{-1.5em}

\begin{abstract}
本项目研究 Text-to-SQL 系统中的候选选择问题：在上游模型已经生成若干条 SQL 后，如何从中选出真正符合问题语义的一条。项目最初计划训练新的 Schema Linker，提高表和列的预测准确率；在分析候选池后，研究重点转向了一个更容易独立验证的问题，即保持原 Schema Linker 和候选生成器不变，训练监督式 Schema-aware reranker。当前已经完成 BIRD 数据处理、固定候选池构建、候选执行标注、结构与 schema 特征提取、GPU pairwise MLP、冻结 MiniLM 语义表示、残差融合和置信门控。现有50题测试中，自一致性准确率为62\%，候选池上界为70\%；当前模型最高只能追平基线，schema-aware 版本最高为60\%。这个结果没有支持最初的性能假设，但揭示了更具体的问题：有效正负候选对过少，简单 schema 特征在同题候选之间缺乏差异，而且验证与执行标签仍有噪声。下一阶段将优先修复数据和评估协议，再训练真正面向 SQL 语义差异的 reranker。
\end{abstract}

\section{问题从哪里来}

Text-to-SQL 系统通常不只生成一条答案。通过不同提示、采样或搜索路径，可以得到多条语法正确且能够执行的 SQL。常见做法是按执行结果聚类，选择最大结果簇中的候选。这种自一致性方法简单而稳定，但它依赖一个并不总成立的假设：多数候选对应的结果更可能正确。如果生成器对问题产生了同一种误解，例如选错相似列、遗漏过滤条件或采用错误连接，多条错误 SQL 也会形成最大的结果簇。

项目开始时，我更关注上游 Schema Linking，计划用神经网络预测问题涉及的表和列。然而，这会同时改变候选的生成分布，很难判断最终增益究竟来自 linker、生成器还是选择器。对固定候选池的检查给出了更直接的切入点：50个开发集问题上，自一致性只选对31题，而候选池中有35题至少包含一条正确 SQL。也就是说，现有候选中已经有答案，但选择规则浪费了4题、8个百分点的空间。

因此，项目最终聚焦于一个独立问题：给定问题、数据库 schema 和固定候选集合，训练一个 selector 对候选重新排序。这个定位使候选生成能力和候选选择能力可以分别测量，也避免把项目结果归因于上游搜索预算的变化。

\section{文献带来的主要启发}

早期工作主要帮助我理解 schema 信息应该如何进入模型。SLSQL 的分析表明，Schema Linking 对跨数据库泛化非常关键；在提供高质量 linking 时，一个相对简单的生成模型也能得到较强结果\cite{lei2020slsql}。这促使我保留原 linker 输出，但由于监督式 linker 本身也并不完美，它不能成为硬过滤条件。

BRIDGE 将问题、schema 和问题中出现的数据库值放进统一序列，通过 BERT 学习文本与表格结构之间的依赖\cite{lin2020bridge}。RAT-SQL 则进一步把问题到表列的匹配、表列包含关系以及主外键关系编码为显式关系\cite{wang2020ratsql}。这两项工作说明，schema-aware 不应只等价于字符串重合率；字段语义、关系图和问题上下文需要联合建模。当前实现中的 candidate-used tables、columns、foreign-key connectivity 和 join completeness 都来自这部分启发。

较新的工作使研究重点转向候选选择。MSc-SQL 让 critic 同时观察多条候选及其执行元数据，在固定候选池上将 BIRD Dev 的结果从 self-consistency 的62.6\%提高到65.6\%，而候选池 oracle 为71.4\%\cite{gorti2025mscsql}。这说明 selector 本身确实有可优化空间。CHASE-SQL 使用经过训练的二候选比较模型，将所有候选两两比较并累计胜出分数\cite{pourreza2024chasesql}，直接启发了本项目的 pairwise RankNet 目标。DPC 则从另一个方向指出，模型只看 SQL 文本时容易受到确认偏差和符号推理能力限制；它通过最小区分数据库和 SQL/Python 双范式验证来暴露候选间的逻辑差异\cite{li2026dpc}。这一思路尚未实现，但为后续构造更可靠的难例和执行验证提供了方向。

\section{最终研究方案}

本项目不重新训练 Schema Linker，也不要求候选来自某一种特定搜索算法。上游只负责提供固定候选池，selector 负责回答“哪一条更值得选”。整体流程如下：

\begin{center}
\small
\begin{tabular}{c}
问题、evidence 与数据库 schema\\
$\downarrow$\\
原 Schema Linker 输出 $+$ 固定候选 SQL 集合\\
$\downarrow$\\
SQL/AST、候选使用的表列、外键路径与执行反馈\\
$\downarrow$\\
Pairwise Schema-aware reranker\\
$\downarrow$\\
自一致性残差融合或置信门控\\
$\downarrow$\\
最终 SQL
\end{tabular}
\end{center}

候选表示分为四部分。第一部分是 SQL 结构，包括表列数量、谓词、聚合、分组、排序、子查询和解析状态。第二部分是 schema 信息，包括候选实际使用的表列、与 linker 输出的覆盖关系、未链接元素、外键连通性和 join key。第三部分是执行信息，包括执行成功、结果行列数、空值比例和执行结果簇大小。第四部分是可选的生成轨迹信息，例如路径深度和已有一致性分数。

对于同一问题的一条正确候选 $s^+$ 和错误候选 $s^-$，模型学习
\[
f(q,s^+) > f(q,s^-),
\]
并使用 RankNet 损失
\[
\mathcal{L}_{\mathrm{rank}}
=-\log \sigma\!\left(f(q,s^+)-f(q,s^-)\right).
\]
推理时不必无条件替换自一致性。当前实现可以将神经分数作为 residual 加到结果簇分数上，也可以只在模型相对基线候选具有足够大置信差时改票。这一设计的目的，是在保留基线稳定性的同时尝试恢复少数 oracle-gap 问题。

\section{已经完成的实现}

项目已经形成从数据到评估的完整离线流程。BIRD train/dev 被安装到本地后，脚本可以按数据库采样问题、生成候选、保存中间结果，并在中断后继续运行。为了尽快获得足够多的训练问题，后期使用直接候选生成绕开了完整 MCTS；这不会改变研究目标，因为所有 selector 都在相同的离线候选池上比较。

目前实现的模型包括 Logistic Regression、小型 GPU MLP 和冻结的 MiniLM 编码器。数值模型的 SQL 加执行版本有25个特征，加入 schema 后为40个，再加入 MCTS 信息后为44个。语义版本将 ``question + evidence'' 与 ``candidate SQL + candidate schema + linked schema'' 组成文本对，并将冻结编码器的表示与数值特征拼接。模型训练、checkpoint、逐候选分数和汇总报告均已保存。

\begin{table}[ht]
\centering
\caption{当前候选池的实际规模。}
\begin{tabular}{lr}
\toprule
项目 & 数量\\
\midrule
请求生成的训练问题 & 300\\
成功保存候选的问题 & 286\\
可执行 gold 并进入训练统计的问题 & 269\\
真实训练候选 & 903\\
至少包含一条正确候选的问题 & 213\\
同时包含正确和错误候选的问题 & 40\\
测试问题 / 测试候选 & 50 / 186\\
\bottomrule
\end{tabular}
\end{table}

这里最值得注意的是最后第三行。903条候选看起来并不少，但真正能够形成同题正负排序监督的只有40个问题。大量问题的候选要么全部执行正确，要么全部错误，因此不能训练“同一个问题中哪一条更好”。报告中的实际训练 pair 只有约111--115个。

\section{实验结果}

所有方法在同一组50题、186条候选上评估。Self-consistency 选对31题，候选池 oracle 为35题。换言之，15题根本没有正确候选，31题已经被基线选对，真正可能由 selector 恢复的只有4题。

\begin{table}[ht]
\centering
\caption{固定测试候选池上的主要结果。}
\begin{tabularx}{0.88\textwidth}{Xcc}
\toprule
方法 & EX & 相对基线\\
\midrule
Execution self-consistency & \textbf{62\%} & --\\
Candidate oracle & 70\% & $+8$\\
数值 MLP，SQL + execution & 60\% & $-2$\\
数值 MLP，schema-aware & 58--60\% & $-2$ 至 $-4$\\
冻结 MiniLM，SQL + execution & 60--62\% & $0$ 至 $-2$\\
冻结 MiniLM，schema-aware & 58--60\% & $-2$ 至 $-4$\\
Schema + execution + MCTS & 58--60\% & $-2$ 至 $-4$\\
\bottomrule
\end{tabularx}
\end{table}

目前没有模型超过62\%的基线。最好的 SQL-only 模型追平基线，最好的 schema-aware 模型为60\%。在一个代表性的 schema-aware 结果中，模型成功修正了2题，但把原本正确的3题改错，因此净结果从31/50下降为30/50。这个结果说明模型并非完全不会识别错误候选，但它的置信度还不足以支持稳定改票。

\section{结果如何理解}

第一个问题是监督数据的有效规模。Pairwise reranker 需要同一问题中的正负候选，而当前269个问题中只有40个满足条件。继续增加“全部错误”或“全部正确”的问题，对排序训练帮助很小。下一轮数据构建应该把目标从候选总数改为正负共存问题数。

第二个问题是当前 schema 特征缺乏区分度。在40个正负混合训练问题中，正确和错误候选的 table coverage 在36题中完全相同，column coverage 在26题中相同，foreign-key connectivity 在39题中相同。很多错误 SQL 与正确 SQL 使用相同表列，只在过滤、聚合或排序上不同。由此可见，简单的 schema overlap 无法解决候选选择问题。

第三个问题是错误候选也可能非常符合 linker。候选通常由相同问题和相同 linker 输出生成，系统性错误会沿着同一 schema 传播。因此，``与 linker 一致''不能直接解释为``语义正确''。真正有用的 schema-aware 模型需要理解问题短语与候选操作之间的对应关系，例如“最高”是否落实为正确的排序和限制，“平均”作用于哪一列，以及过滤条件是否落在正确实体上。

第四个问题是当前所谓的语义模型仍然较弱。MiniLM 原本针对文本检索训练，而且在本实验中完全冻结。它提供的是通用相关性表示，并没有针对 SQL 运算、聚合、连接和执行状态进行训练。因此，现有结果只能说明冻结的通用语义特征没有帮助，不能据此否定候选级语义建模。

最后，实验协议还有两处需要修正。Gold-positive 消融目前在划分训练和验证集之前注入 gold，导致验证集中也出现人工候选，内部验证分数因此明显虚高。执行标签还会无条件忽略结果行顺序，可能把错误的 ORDER BY 候选标为正确。这两项问题都应在扩大实验前修复。

\section{下一阶段计划}

下一阶段不会首先扩大模型，而是按以下顺序推进：

\begin{enumerate}
  \item \textbf{修复评估协议。} 先按数据库划分，再只向训练部分注入 gold；接入官方 BIRD evaluator，并正确处理 ORDER BY。
  \item \textbf{扩大有效排序数据。} 目标是至少500个正负共存问题，并定向构造错误列、错误 join、漏条件、错误聚合和错误排序等可执行难例。
  \item \textbf{训练真正的 SQL 语义模型。} 输入问题、evidence、候选 AST、候选使用的具体表列描述和外键路径，比较冻结、部分解冻和完整微调。
  \item \textbf{建模候选之间的差异。} 从独立打分发展到 pairwise tournament 或 groupwise critic，显式指出两条 SQL 在谓词、聚合、连接和排序上的差别。
  \item \textbf{采用保守决策。} 在数据库隔离验证集上校准门控阈值，主要报告正确改票数、错误改票数和净收益，并使用至少五个随机种子。
\end{enumerate}

\section{阶段性结论}

项目已经完成候选池构建、执行标注、schema 特征、GPU pairwise 训练、语义表示和选择策略的端到端离线流程，但当前结果没有证明 schema-aware reranker 能提高准确率。这个负结果仍然有价值，因为它排除了一个过于简单的假设：只要把 schema overlap 和冻结语言模型表示加进 MLP，selector 就会自然变好。

目前更合理的判断是，项目的首要瓶颈不是算力，而是有效同题正负对不足，以及候选级语义差异没有被当前特征表示出来。后续工作的重点将从“加入更多特征”转向“构造更有判别力的数据，并直接学习问题与候选逻辑之间的对应关系”。如果经过严格数据库隔离和多随机种子验证后，schema-aware 模型仍不能超过 SQL + execution 模型，那么项目也能够给出一个清晰的研究结论：原 Schema Linker 的硬选择结果对候选重排序没有提供稳定的额外信息。

\begin{thebibliography}{9}
\small

\bibitem{lin2020bridge}
Xi Victoria Lin, Richard Socher, and Caiming Xiong.
\newblock Bridging Textual and Tabular Data for Cross-Domain Text-to-SQL Semantic Parsing.
\newblock \emph{Findings of EMNLP}, 2020.

\bibitem{lei2020slsql}
Wenqiang Lei, Weixin Wang, Zhixin Ma, Tian Gan, Wei Lu, Min-Yen Kan, and Tat-Seng Chua.
\newblock Re-examining the Role of Schema Linking in Text-to-SQL.
\newblock \emph{EMNLP}, 2020.

\bibitem{wang2020ratsql}
Bailin Wang, Richard Shin, Xiaodong Liu, Oleksandr Polozov, and Matthew Richardson.
\newblock RAT-SQL: Relation-Aware Schema Encoding and Linking for Text-to-SQL Parsers.
\newblock \emph{ACL}, 2020.

\bibitem{gorti2025mscsql}
Satya Krishna Gorti et al.
\newblock MSc-SQL: Multi-Sample Critiquing Small Language Models for Text-to-SQL Translation.
\newblock \emph{NAACL}, 2025.

\bibitem{pourreza2024chasesql}
Mohammadreza Pourreza et al.
\newblock CHASE-SQL: Multi-Path Reasoning and Preference Optimized Candidate Selection in Text-to-SQL.
\newblock \emph{arXiv:2410.01943}, 2024.

\bibitem{li2026dpc}
Boyan Li, Ou Ocean Kun Hei, Yue Yu, and Yuyu Luo.
\newblock DPC: Training-Free Text-to-SQL Candidate Selection via Dual-Paradigm Consistency.
\newblock \emph{ACL}, 2026.

\end{thebibliography}

\end{document}
